% Generated by results/thesis_tables.py - do not edit by hand.
\begin{table}[H]
\centering
\footnotesize
\setlength{\tabcolsep}{3pt}
\begin{tabular}{|l|r|r|r|r|>{\centering\arraybackslash}p{1.35cm}|>{\centering\arraybackslash}p{1.35cm}|>{\centering\arraybackslash}p{1.35cm}|>{\centering\arraybackslash}p{1.35cm}|}
\hline
Chains & Κακόβ. & Νόμιμα & Ομοιότ. (κακόβ.) & Ομοιότ. (νόμιμα) & Πλήρες σύστημα & Μόνο δείγματα & Μόνο πολιτικές & Χωρίς RAG \\
\hline
1--10 & 299 & 206 & 0,965 & 0,937 & 290/299 / 3 & 287/299 / 1 & 288/299 / 3 & 285/299 / 1 \\
11--100 & 127 & 122 & 1,000 & 0,957 & 125/127 / 2 & 125/127 / 4 & 123/127 / 2 & 123/127 / 3 \\
$>100$ & 16 & 15 & 0,990 & 0,976 & 13/16 / 0 & 14/16 / 0 & 14/16 / 1 & 14/16 / 0 \\
\hline
\end{tabular}
\caption{Αποτελέσματα ανά πλήθος Behavior Chains του πακέτου: διάμεση μέγιστη ομοιότητα με το ευρετήριο και, σε κάθε διαμόρφωση, κακόβουλα πακέτα που εντοπίστηκαν / ψευδώς θετικά. Συσχέτιση Spearman ανάμεσα στο πλήθος των Behavior Chains και στη μέγιστη ομοιότητα: 0,54 στα κακόβουλα και 0,44 στα νόμιμα πακέτα.}
\label{tab:chain_buckets}
\end{table}
